
\needspace{4\baselineskip}
\color{uniblue}{\subsection{\sloppy Логика отключения при срабатывании технологической сигнализации \mbox{трансформатора} (ЛО~ТС)}}
\color{black}

\begin{enumerate}

%%%%%%%%%%%%%%%%%%%%%%%%%%%%%%%%%%%%%%%%%%%%%%%%%%%%%%%%%%%%%%  ОБЩИЕ ПАРАМЕТРЫ  %%%%%%%%%%%%%%%%%%%%%%%%%%%%%%%%%%%%%%%%%%%%%%%%%%
\item Общие сведения

\begin{enumerate}

\item
Логика отключения при срабатывании технологической сигнализации служит для реализации действия на отключение трансформатора при срабатывании датчиков технологической сигнализации трансформатора в случае необходимости или требований.

\item
В устройстве логика отключения при срабатывании технологической сигнализации представлена функциональным блоком «ЛО ТС». В функциональном блоке реализована возможность приема следующих сигналов датчиков технологической сигнализации трансформатора:
\begin{itemize}
\item открытое состояние предохранительного клапана для ЛО ПК (логика отключения при срабатывании предохранительного клапана);
\item закрытое состояние отсечного клапана для функции ЛО ОК (логика отключения при срабатывании отсечного клапана);
\item низкий уровень масла в расширительном баке трансформатора для ЛО ДУм расширителя (логика отключения при снижении уровня масла в расширительном баке трансформатора).
\end{itemize}

\item
ЛО~ТС может быть выведена из работы оперативно путем активации сигнала «ОВ~ЛО~ТС».
\item
В таблице \ref{lotspk:tbl1} приведены параметры, необходимые для настройки функций в составе ЛО ТС.

\end{enumerate}
%%%%%%%%%%%%%%%%%%%%%%%%%%%%%%%%%%%%%%%%%%%%%%%%%%%%%%%%%%%%%%%%%%%%%%  ПРЕДОХРАНИТЕЛЬНЫЙ КЛАПАН  %%%%%%%%%%%%%%%%%%%%%%%%%%%%%%%%%%%%%%%%%%%%%%%%%%

\item Логика отключения при срабатывании предохранительного клапана (ЛО~ПК)

\begin{enumerate}

\item
Логика отключения при срабатывании предохранительного клапана (ЛО ПК) служит для реализации действия на отключение трансформатора при срабатывании предохранительного клапана трансформатора.

\item
В функциональном блоке «ЛО ТС» логика отключения при срабатывании предохранительного клапана представлена функцией «ЛО ПК».

\item
Ввод функции в работу на этапе параметрирования устройства осуществляется программным переключателем SGF1 «Ввод\_функции». Информация о введенной функции отображается выходным сигналом «ЛО ПК: Ввод».

\item
Функция может быть выведена из работы оперативно путем активации сигнала «ОВ ЛО ПК», что характеризуется активным состоянием сигнала «ЛО ПК: Оперативный вывод» на выходе функции.

\item
При активном состоянии входного сигнала «ПК~на~сигн» действие отключающей логики переводится на сигнал.

\item
Алгоритм работы функции выполнен в соответствии с рисунком \ref{lotspk:img1}. 

\vspace{1mm}
\begin{figure}[H]
\centering
\includegraphics[width=1\textwidth,height=1\textheight,keepaspectratio]{\fbpath/20. ЛО ТС Т/_latex/img/lotspk.pdf}
\captionof{figure}{Функциональная схема ЛО ПК}
\label{lotspk:img1}
\end{figure}

\end{enumerate}
\FloatBarrier % Все float'ы ДО этой строки будут размещены здесь.
%%%%%%%%%%%%%%%%%%%%%%%%%%%%%%%%%%%%%%%%%%%%%%%%%%%%%%%%%%%%%%%%%%%%%% ОТСЕЧНОЙ КЛАПАН %%%%%%%%%%%%%%%%%%%%%%%%%%%%%%%%%%%%%%%%%%%%%%%%%%
\item Логика отключения при срабатывании отсечного клапана (ЛО~ОК)

\begin{enumerate}

\item
Логика отключения при срабатывании отсечного клапана (ЛО~ОК) служит для реализации действия на отключение трансформатора при срабатывании отсечного клапана трансформатора.

\item
В функциональном блоке «ЛО ТС» логика отключения при срабатывании отсечного клапана представлена функцией «ЛО ОК».

\item
Ввод функции в работу на этапе параметрирования устройства осуществляется программным переключателем SGF1 «Ввод\_функции». Информация о введенной функции отображается выходным сигналом «ЛО ОК: Ввод».

\item
Функция «ЛО ПК» может быть выведена из работы оперативно путем активации сигнала «ОВ ЛО ОК», что характеризуется активным состоянием сигнала «ЛО ОК: Оперативный вывод» на выходе функции.

\item
При активном состоянии входного сигнала «ОК~на~сигн» действие отключающей логики переводится на сигнал.

\item
Алгоритм работы функции выполнен в соответствии с рисунком \ref{lotsok:img1}.

\vspace{1mm}
\begin{figure}[H]
\centering
\includegraphics[width=1\textwidth,height=1\textheight,keepaspectratio]{\fbpath/20. ЛО ТС Т/_latex/img/lotsok.pdf}
\captionof{figure}{Функциональная схема ЛО ОК}\label{lotsok:img1}
\end{figure}

\end{enumerate}
\FloatBarrier % Все float'ы ДО этой строки будут размещены здесь.
%%%%%%%%%%%%%%%%%%%%%%%%%%%%%%%%%%%%%%%%%%%%%%%%%%%%%%%%%%%%%%%%%%%%%%  НИЗКИЙ УРОВЕНЬ МАСЛА В РАСШИРИТЕЛЕ %%%%%%%%%%%%%%%%%%%%%%%%%%%%%%%%%%%%%%%%%%%%%%

\item Логика отключения при снижении уровня масла в расширительном баке трансформатора (ЛО~ДУм расширителя)

\begin{enumerate}

\item
Логика отключения при снижении уровня масла в расширительном баке трансформатора (ЛО~ДУм расширителя) служит для реализации действия на отключение трансформатора при снижении уровня масла в расширительном баке трансформатора.

\item
В функциональном блоке «ЛО ТС» логика отключения при снижении уровня масла в расширительном баке трансформатора представлена функцией «ЛО~ДУм расширителя».

\item
Ввод функции в работу на этапе параметрирования устройства осуществляется программным переключателем SGF1 «Ввод\_функции». Информация о введенной функции отображается выходным сигналом «ЛО~ДУм расширителя: Ввод».

\item
Функция может быть выведена из работы оперативно путем активации сигнала «ОВ ЛО ДУм расшир», что характеризуется активным состоянием сигнала «ЛО~ДУм расшир: Оперативный вывод» на выходе функции.

\item
При активном состоянии входного сигнала «ДУм расшир~на~сигн» действие отключающей логики функции переводится на сигнал. 

\item
Алгоритм работы функции выполнен в соответствии с рисунком \ref{lotsdut:img1}. 

\vspace{1mm}
\begin{figure}[H]
\centering
\includegraphics[width=1\textwidth,height=1\textheight,keepaspectratio]{\fbpath/20. ЛО ТС Т/_latex/img/lotsdut.pdf}
\captionof{figure}{Функциональная схема ЛО~ДУм расширителя}\label{lotsdut:img1}
\end{figure}

\end{enumerate}

\end{enumerate}
\FloatBarrier % Все float'ы ДО этой строки будут размещены здесь.